\section{Who can bind the state that enforces the rule?}
\label{sec:13.7}
A state can bind a private developer because it controls something the developer needs. When the state is the operator, the institution that writes the rule, controls the evidence and supplies the remedy is the institution the rule is meant to constrain.

The chapter's cases are American, and “the state” is tested here on one constitutional order. Europe is a different one: the Convention applies, has extraterritorial reach, and produces judgments states sometimes obey; there is no state-secrets privilege of the American kind and no collapsed damages remedy. The AI Act makes the boundary visible: it can require documentation from private providers, subject high-risk systems to review, and impose penalties tied to global turnover — and it excludes military, defence and national-security uses from its scope. Public authorities stay covered elsewhere, but enforcement against a member state's agencies still runs through institutions that state constitutes.

For review to bind here, the chain that makes review binding has to run against a state: initiation by people outside its electorate, evidence it would withhold, an interim order before the use is complete, and a remedy that doesn't need its consent. A principle supplying none of those may state the rule correctly and identifies nobody who can make the state obey it.

A treaty obligation of that kind exists: the Council of Europe's Framework Convention, binding at the level of a human-rights instrument with operational detail left to domestic law \autocite{coe2024framework} — and not yet in force, needing five ratifications, of which the EU deposited one in May 2026 \autocite{coe2026euratifies}. Everything else at that level is voluntary.

Every instrument I propose binds through one of three channels: a licence term, an appropriation condition, or a statute. The first two run through procurement, and procurement is one office. A captured Congress waives the appropriation condition; one memorandum from the Office of Management and Budget strikes the licence terms from every agreement at once. Three second routes exist and none is as strong. A collective agreement carries codetermination and the individual refusal right without the state. A coalition partner's precondition carries the review ratio into any war the partner takes part in. State and municipal procurement carry the civilian terms into deployments the federal government doesn't run. The force term has no second route, which is why it is a statute.

\textbf{Residual risk.} One residual risk no design choice removes, and it's what the floor was introduced for. Every lever belongs to somebody, and a government that can reach whoever holds the lever has the lever. A constraint that holds against the one in possession has to be held by something it cannot instruct, and under total capture nothing holds. Every instrument in this book ends at this limit; it is stated here and not again.
